\chapter{Positioning and Sentiment}\label{ch:s1:positioning-and-sentiment}

Every Friday at 3:30 in the afternoon, New York time, the CFTC publishes who held what in each futures market on the Tuesday before: producers and
merchants, swap dealers, managed money. The report is public, cheap and old by the time it appears, and it has been studied for as long as it has
existed. De Roon, Nijman and Veld found that hedgers' positions help explain futures returns across twenty markets: speculators who take the other side
of hedging demand are paid for it. On this chapter's synthetic commodities, where that payment is planted, the hedgers' position has a rank IC of 0.044
with the next week's returns when it is known at once, 0.038 when it is known on Friday, and 0.017 three weeks later. On nine years of real corn
positioning, the most studied signal in the market is not significant. The build is \texttt{firm.posisig}.

\section{Positioning reports}

\begin{dated}{2026-09}{The Commitments of Traders schedule}\label{dat:s1:positioning-and-sentiment:cot}
The CFTC's Commitments of Traders report is generally published each Friday at 3:30 pm Eastern Time with positions as of the immediately preceding
Tuesday; the Commission says it takes three days to process the data, and holidays shift the schedule, which it publishes in advance. Book~3,
chapter~9 describes the report's trader categories.
\end{dated}

\begin{definition}[Hedging pressure, positioning signal]\label{def:s1:positioning-and-sentiment:hedging}
\emph{Hedging pressure}\index{hedging pressure} is the net position of hedgers in a futures market, usually net short for producers, as a share of
open interest; speculators who take the other side earn a risk premium that rises with it. A \emph{positioning signal}\index{positioning signal} is a
trading signal built from reported positions of a category of traders, as known at the time of the trade.
\end{definition}

The record is older than the chapter's other signals. Bessembinder found that returns in currency and agricultural futures vary with hedgers' net
holdings after controlling for systematic risk; de Roon and co-authors found that both a market's own \omterm{def:s1:positioning-and-sentiment:hedging}{hedging pressure} and that of related markets in
its group affect returns, and that the effect survives a control for \omterm{def:s1:flows:fit}{price pressure}. The logic is the insurance view of futures: hedgers pay to lay off
risk, and whoever carries it is paid.

In real positions the two sides are mirror images. In the CFTC's disaggregated data for CBOT corn, 2016 to 2026, managed money's net position ranged
from $-22.8\%$ to $+24.4\%$ of open interest, and its correlation with the producers' and merchants' net position is $-0.97$
(\cref{fig:s1:positioning-and-sentiment:corn}): speculators absorb what hedgers lay off.

\begin{omfigure}
\begin{tikzpicture}
\begin{axis}[width=11cm,height=5.2cm,xlabel={\small year},ylabel={\small net position (\% of open interest)},tick label style={font=\footnotesize},
  xmin=2016,xmax=2026.8,ymin=-55,ymax=30,grid=major,grid style={gray!25},legend columns=2,legend style={font=\scriptsize,at={(0.5,-0.3)},
  anchor=north,draw=none,fill=none,/tikz/every even column/.append style={column sep=8pt}},table/col sep=comma,/pgf/number format/1000 sep={}]
\addplot[omThm,thick] table[x=year,y=mm]{figdata/strategies-1/24-positioning-and-sentiment/corn.csv};
\addplot[omDef,thick] table[x=year,y=pm]{figdata/strategies-1/24-positioning-and-sentiment/corn.csv};
\addplot[black!40,dashed] coordinates {(2016,0) (2026.8,0)};
\legend{managed money,producers and merchants}
\end{axis}
\end{tikzpicture}

\omcaption{CBOT corn futures: net positions of managed money and of producers, merchants, processors and users as a share of open interest, weekly
(Tuesday), January 2016 to September 2026, from the CFTC's disaggregated Commitments of Traders. Data: \texttt{s1\_posisig.corn}.}\label{fig:s1:positioning-and-sentiment:corn}
\end{omfigure}

Does corn's managed money position predict corn's price? The test uses the position's 52-week z-score as published by each month's end and the IMF
maize price's change over the month after next (monthly averages; skipping a month keeps the signal's own month out of the return). Over 113 months,
2017 to 2026, the correlation is 0.098 ($t = 1.04$). After record-long positions (z-score above 1) the next three months' price change averaged
$+4.6\%$ (28 cases), after record-short ones $+0.2\%$ (24), and $-0.9\%$ otherwise (59): if anything, speculators' extremes were followed by moves in
their favour, as the hedging-pressure view would have it, but with a hundred months the evidence says nothing either way.

\section{The value of a report, and its age}

\texttt{firm.posisig} plants the mechanism in \texttt{firm.synthfut}'s ten commodities (\cref{lst:s1:positioning-and-sentiment:model}). Hedgers' net
short position follows an AR(1) with a half-life of 20 trading days, in standard deviations; it adds a premium of 0.3 times the market's volatility a
year per standard deviation to the next day's return. Speculators' reported net long position is the \omterm{def:s1:positioning-and-sentiment:hedging}{hedging pressure} plus half a standard deviation
per unit of their own 60-day trend, plus noise. Positions are measured on the fifth day of each week (the ``Tuesday'') and known after a publication
lag. The signal is measured every day against the next five days' returns, across the ten markets.

\begin{center}
\small
\begin{tabular}{@{}lcccc@{}}
\toprule
publication lag (trading days) & 0 & 3 (Friday) & 10 & 21\\
\midrule
hedgers' position: rank IC & 0.044 & 0.038 & 0.029 & 0.017\\
\quad t-statistic & 5.1 & 4.4 & 3.4 & 2.0\\
speculators' position: rank IC & 0.043 & 0.039 & 0.031 & 0.024\\
\bottomrule
\end{tabular}
\end{center}

The Tuesday-to-Friday lag costs 14\% of the signal (\cref{fig:s1:positioning-and-sentiment:lags}); three weeks cost 62\%. Both follow from the
pressure's persistence: with a half-life of 20 days, three days of age keep 90\% of it, 21 days 48\%. The speculators' position loses its value more
slowly because part of it is trend-chasing, and \texttt{firm.synthfut}'s planted trends are slow: the trend component of the signal ages well even
though it has nothing to do with hedging.

\begin{omfigure}
\begin{tikzpicture}
\begin{axis}[width=10cm,height=5cm,xlabel={\small publication lag (trading days)},ylabel={\small rank IC, next five days},
  tick label style={font=\footnotesize},xmin=-1,xmax=22,ymin=0,ymax=0.05,grid=major,grid style={gray!25},legend pos=north east,
  legend style={font=\scriptsize},scaled y ticks=false,yticklabel style={/pgf/number format/fixed,/pgf/number format/precision=2},
  table/col sep=comma]
\addplot[omThm,thick,mark=*] table[x=lag,y=hedge]{figdata/strategies-1/24-positioning-and-sentiment/lags.csv};
\addplot[omDef,thick,mark=square*] table[x=lag,y=spec]{figdata/strategies-1/24-positioning-and-sentiment/lags.csv};
\legend{hedgers' net short,speculators' net long}
\end{axis}
\end{tikzpicture}

\omcaption{\omterm{def:s1:positioning-and-sentiment:hedging}{Positioning signals} on the synthetic commodities with planted \omterm{def:s1:positioning-and-sentiment:hedging}{hedging pressure} (half-life 20 days): the rank IC with the next five days' returns
of the hedgers' and the speculators' reported positions, by the lag between the measurement and its publication. Data:
\texttt{s1\_posisig.lag\_table}.}\label{fig:s1:positioning-and-sentiment:lags}
\end{omfigure}

The contrarian rule, selling when speculators are at a record long, fails in the model by construction: after a 52-week z-score above 1.5 the next
quarter's return averaged $+0.045$ annual volatilities, after one below $-1.5$ it averaged $-0.093$, and $-0.034$ otherwise. A contrarian rule needs a
mechanism in which speculators are the ones paying, crowding into positions that later unwind; the model has none, and neither the corn data nor the
sources fetched for this chapter supply one.

\section{Surveys and sentiment indices}

\begin{definition}[Sentiment index]\label{def:s1:positioning-and-sentiment:sentiment}
A \emph{sentiment index}\index{sentiment index} summarises investors' optimism or pessimism from surveys or market-based proxies (issuance, fund
flows, first-day returns of new issues, valuation spreads), usually as the first principal component of several proxies.
\end{definition}

Baker and Wurgler showed what sentiment predicts in the cross-section of stocks: when sentiment is low at the start of a period, the stocks hardest to
value and to arbitrage (small, young, volatile, unprofitable, non-dividend-paying, extreme growth and distressed stocks) earn relatively high returns
later; when it is high, they earn relatively low ones. Sentiment is a conditioning variable more than a signal: it says which stocks a wave of
optimism has moved furthest from value, and when their mispricing is likely to be large. It changes slowly and is published with a lag, so its age
matters less than a weekly positioning report's.

Options positioning is the third source: \omterm{def:s1:options-implied-signals-for-stocks:volume}{put--call ratios}, open interest by strike and dealers' likely gamma. Chapter~15 measured what informed option
demand says about stocks; chapter~13 measured what dealers' hedging does near the close. As positioning, options add the view of the traders who pay
for convexity, reported daily and by strike.

\section{Strategy files}

\begin{strategyfile}{Hedging-pressure signal}\label{strat:s1:positioning-and-sentiment:hedging}
\sfield{Who pays you, and why} Hedgers who pay a premium to lay off price risk.
\sfield{Instruments and venues} Commodity and financial futures with COT reporting.
\sfield{Signal} Hedgers' net short position as a share of open interest, as published (Friday), within market or across a group.
\sfield{Sizing and execution} Long markets with high \omterm{def:s1:positioning-and-sentiment:hedging}{hedging pressure}, short low; weekly, volatility-scaled.
\sfield{Costs} Low; weekly rebalancing.
\sfield{How it dies} More speculative capital competing to absorb hedging; changing hedging programmes.
\sfield{Horizon, capacity, infrastructure} Weeks; the COT archive with release dates.
\sfield{Backtest honestly} Positions only from the Friday of each report (later after holidays), never from Tuesday.
\sfield{Sources} Bessembinder (1992); de Roon, Nijman and Veld (2000); this chapter: IC 0.038 as published, synthetic.
\end{strategyfile}

\begin{strategyfile}{Extreme-positioning contrarian}\label{strat:s1:positioning-and-sentiment:extreme}
\sfield{Who pays you, and why} Only if speculators crowd into positions and unwind them together; otherwise nobody.
\sfield{Instruments and venues} Futures with COT reporting.
\sfield{Signal} Managed money's net position at a multi-year extreme.
\sfield{Sizing and execution} Fade the extreme, small, with a stop.
\sfield{Costs} Low.
\sfield{How it dies} It may not live: in corn, 2017--2026, extremes were followed by moves in the speculators' favour.
\sfield{Horizon, capacity, infrastructure} Months.
\sfield{Backtest honestly} Extremes as measured with past data only; published dates; enough extremes to test.
\sfield{Sources} No supporting result verified; this chapter's corn test ($t = 1.04$ for the opposite sign).
\end{strategyfile}

\begin{strategyfile}{Options-positioning signal}\label{strat:s1:positioning-and-sentiment:options}
\sfield{Who pays you, and why} Traders who pay for protection or leverage, whose demand is visible in option volumes and open interest.
\sfield{Instruments and venues} Stocks and indices with listed options.
\sfield{Signal} \omterm{def:s1:options-implied-signals-for-stocks:volume}{Put--call ratios} and open interest by strike, conditioned on sentiment.
\sfield{Sizing and execution} Tilts in stock or index positions; daily.
\sfield{Costs} Moderate.
\sfield{How it dies} Hedging flow that looks like speculation; crowding.
\sfield{Horizon, capacity, infrastructure} Days to weeks; options data by strike.
\sfield{Backtest honestly} Open interest as reported the next morning; volumes signed only if the data allow.
\sfield{Sources} Chapter~15; Baker and Wurgler (2006) for sentiment as a conditioning variable.
\end{strategyfile}

\section{Tutorial: record long}\label{tut:s1:positioning-and-sentiment:tut}

\begin{tutorial}
\textbf{Goal.} Plant \omterm{def:s1:positioning-and-sentiment:hedging}{hedging pressure} in synthetic commodities, measure what the \omterm{def:s1:positioning-and-sentiment:hedging}{positioning signal} is worth at each publication lag and at extremes,
and test the real corn positioning series. \textbf{End state:} the table and the two figures.
\begin{tutsteps}
\item \textbf{Reports and the market}: the publication schedule and the planted \omterm{def:s1:positioning-and-sentiment:hedging}{hedging pressure}.
\omcode{code/firm/posisig/firm_posisig.py}{26}{58}{The reporting schedule and a market with hedging pressure.}\label{lst:s1:positioning-and-sentiment:model}
\item \textbf{Real corn}: managed money's z-score as published by month-end against maize prices.
\omcode{code/strategies-1/24-positioning-and-sentiment/python/s1_posisig.py}{74}{102}{Corn positioning against the maize price.}\label{lst:s1:positioning-and-sentiment:corn}
\item \textbf{Run} \texttt{lag\_table()}, \texttt{extremes()} and \texttt{corn()}, and \texttt{fig\_posisig.py}.
\end{tutsteps}
\textbf{What to change next.} Plant a crowding mechanism (speculators' positions pushing prices, which revert) and see the contrarian rule appear;
shorten the pressure's half-life to five days and measure the Friday lag's cost; add swap dealers as a third category.
\end{tutorial}

\section{Build: positioning signals}\label{bld:s1:positioning-and-sentiment:posisig}

\begin{build}
\bfield{Purpose} Reports as known on each day, a market with \omterm{def:s1:positioning-and-sentiment:hedging}{hedging pressure}, forward returns, rank ICs and z-scores.
\bfield{Interface} \texttt{published(x, every, offset, lag)}, \texttt{hedging\_market(r, vol, half\_life, premium, chase, noise, rng)},
\texttt{forward(r, h)}, \texttt{ic(signal, fwd, step, overlap)}, \texttt{zscore(x, window)}; \texttt{firm.cot}'s \texttt{net\_share} for real reports.
\bfield{Rules} A report is usable only from its publication day; ICs with overlapping horizons have their t-statistics corrected.
\bfield{Acceptance tests} \texttt{code/firm/posisig/tests/}: a Tuesday-to-Friday schedule by hand, forward sums, IC and its overlap correction, the
premium as planted.
\bfield{Stretch} The CFTC's actual release calendar with holiday shifts; cross-market \omterm{def:s1:positioning-and-sentiment:hedging}{hedging pressure} within groups; a crowding mechanism.
\end{build}

\begin{omsources}
\item F.~A.~de Roon, T.~E.~Nijman and C.~Veld, ``Hedging pressure effects in futures markets'', \emph{Journal of Finance} 55(3), 2000.
\item H.~Bessembinder, ``Systematic risk, hedging pressure, and risk premiums in futures markets'', \emph{Review of Financial Studies} 5(4), 1992.
\item M.~Baker and J.~Wurgler, ``Investor sentiment and the cross-section of stock returns'', \emph{Journal of Finance} 61(4), 2006.
\item US Commodity Futures Trading Commission, Commitments of Traders (release schedule and disaggregated reports); IMF, global price of maize (via
FRED).
\end{omsources}

\section{Exercises}

\begin{exercise}[$\star$]\label{exo:s1:positioning-and-sentiment:1}
A position measured on Tuesday is published on Friday. On which day can a backtest first use it, and what does using it on Tuesday do to the results?
\end{exercise}

\begin{exercise}[$\star$]\label{exo:s1:positioning-and-sentiment:2}
\omterm{def:s1:positioning-and-sentiment:hedging}{Hedging pressure} follows an AR(1) with a half-life of 20 trading days. What share of it survives three days? Ten? Twenty-one?
\end{exercise}

\begin{exercise}[$\star$]\label{exo:s1:positioning-and-sentiment:3}
Why are managed money's and producers' positions in corn almost perfect mirror images?
\end{exercise}

\begin{exercise}[$\star\star$]\label{exo:s1:positioning-and-sentiment:4}
With 113 months, how large must a correlation be to be significant at 5\%? Is corn's 0.098?
\end{exercise}

\begin{exercise}[$\star\star$]\label{exo:s1:positioning-and-sentiment:5}
Why does the speculators' signal keep more of its value than the hedgers' at long lags?
\end{exercise}

\begin{exercise}[$\star\star$]\label{exo:s1:positioning-and-sentiment:6}
Why does the contrarian rule fail in the model, and what would have to be planted for it to work?
\end{exercise}

\begin{exercise}[$\star\star\star$]\label{exo:s1:positioning-and-sentiment:7}
\emph{Coding.} Run \texttt{lag\_table(5.0)}, with a half-life of five days instead of twenty. How much of the hedgers' IC does the Friday lag
cost now?
\end{exercise}

\begin{exercise}[$\star\star\star$]\label{exo:s1:positioning-and-sentiment:8}
\emph{Find the flaw.} ``We built a \omterm{def:s1:positioning-and-sentiment:sentiment}{sentiment index} from surveys and it predicted the market's next-month return with $t = 3$ over 1990--2020, using
the survey values as dated.''
\end{exercise}

\section{Problem: Record Long}

\begin{problem}[{Weekend problem --- a report's value and its age}]\label{pb:s1:positioning-and-sentiment:1}
The chapter's synthetic commodities, the corn reports and the public record.

\medskip\noindent\textbf{Part I --- Reports.}
\begin{enumerate}
\item When is the COT report published, and for which day?
\item Define \omterm{def:s1:positioning-and-sentiment:hedging}{hedging pressure} and a \omterm{def:s1:positioning-and-sentiment:hedging}{positioning signal}.
\item What did Bessembinder and de Roon, Nijman and Veld find?
\item How do corn's managed money and producers' positions relate?
\end{enumerate}

\medskip\noindent\textbf{Part II --- Corn.}
\begin{enumerate}[resume]
\item Describe the corn test and why it skips a month.
\item Give its correlation and t-statistic.
\item What followed record longs and record shorts?
\item What does a hundred months allow you to conclude?
\end{enumerate}

\medskip\noindent\textbf{Part III --- The model.}
\begin{enumerate}[resume]
\item Describe the planted \omterm{def:s1:positioning-and-sentiment:hedging}{hedging pressure} and the speculators' position.
\item Give the hedgers' IC at each publication lag.
\item Explain the Friday lag's cost from the half-life.
\item Why does the contrarian rule fail here?
\end{enumerate}

\medskip\noindent\textbf{Part IV --- The verdict.}
\begin{enumerate}[resume]
\item State the \emph{named result}: the IC of the \omterm{def:s1:positioning-and-sentiment:hedging}{positioning signal} and the effect of the Tuesday-to-Friday publication lag.
\item Define a \omterm{def:s1:positioning-and-sentiment:sentiment}{sentiment index}; what did Baker and Wurgler find?
\item How is sentiment used differently from positioning?
\item What does options positioning add?
\item How would you backtest a COT strategy honestly?
\item Which strategy file lacks a mechanism?
\item How does this chapter relate to chapter~20's carry?
\item In one sentence: who pays a positioning trader?
\end{enumerate}
\end{problem}

\section{Interview questions}

\begin{interviewq}[$\star$ \iqroles{researcher}]\label{iq:s1:positioning-and-sentiment:1}
What does the Commitments of Traders report contain, and what can it not tell you?
\end{interviewq}

\begin{interviewq}[$\star\star$ \iqroles{researcher}]\label{iq:s1:positioning-and-sentiment:2}
Why might speculators earn a premium in futures markets?
\end{interviewq}

\begin{interviewq}[$\star\star$ \iqroles{trader}]\label{iq:s1:positioning-and-sentiment:3}
Managed money is at a record long in crude oil. What do you do?
\end{interviewq}

\begin{interviewq}[$\star\star$ \iqroles{developer}]\label{iq:s1:positioning-and-sentiment:4}
How would you store COT data so that backtests use each report only when it was public?
\end{interviewq}

\begin{interviewq}[$\star\star$ \iqroles{researcher}]\label{iq:s1:positioning-and-sentiment:5}
How would you build a \omterm{def:s1:positioning-and-sentiment:sentiment}{sentiment index}, and how would you test it?
\end{interviewq}

\begin{interviewq}[$\star\star\star$ \iqroles{researcher}]\label{iq:s1:positioning-and-sentiment:6}
A signal is an AR(1) with coefficient $\phi$ per day and predicts the next day's return with correlation $\rho$. Derive its correlation with the return
$k$ days after it was measured, and the average correlation for a report used over the five days after a lag of $\ell$ days.
\end{interviewq}
